% GENERATED FILE. Edit canonical lessons, then run npm run generate:books.
% canonical-source-hash: fnv1a64:19acc4006e449b8b
% canonical-lessons: RU-C246-mikrovolnovka, RU-C246-chaynik, RU-C246-miska, RU-C246-blyudtse, RU-C246-kruzhka

\chapter{Microwave, Kettle, Bowl, Saucer, Mug}
\label{ch:ru-microwave-kettle-bowl-saucer-mug}
\hlchaptermodality{246}

\begin{chapteropening}
\textbf{By the end of this chapter:} \emph{I can say a microwave, a kettle, a bowl, a saucer and a mug.}

Complete the last lesson of chapter 246: I can say a microwave, a kettle, a bowl, a saucer and a mug.
\end{chapteropening}

\section[\texorpdfstring{\ru{микроволновка} (mikrovolnóvka)}{микроволновка (mikrovolnóvka)}]{\ru{микроволновка} (mikrovolnóvka) — a microwave}

\label{lesson:RU-C246-mikrovolnovka}



\begin{quote}
Before the new one: say the Russian for terrible, then the Russian for excellent.
\end{quote}

\subsection*{You'll want to know: \ru{микроволновка}}
\textbf{\ru{микроволновка}} (\emph{mikrovolnóvka}) — \textquotedblleft{}a microwave\textquotedblright{}.

\begin{grammarlens}[title={What you've built}]
One more word for this chapter.
\end{grammarlens}

\subsection*{Your turn}
\begin{itemize}
\raggedright
  \item \emph{Say it:} \emph{mikrovolnóvka}
  \item \emph{Say it:} \emph{mikrovolnóvka}, once more
  \item \emph{Say it:} say \emph{mikrovolnóvka}
  \item \emph{Recall:} say the Russian for terrible, then the Russian for excellent, then say \emph{mikrovolnóvka} again
\end{itemize}

\begin{tcolorbox}[breakable,colback=teal!4,colframe=teal!35!black,title={Before you move on}]
What does \ru{микроволновка} mean? (\textquotedblleft{}A microwave\textquotedblright{}.) Say it once more.
\end{tcolorbox}
\section[\texorpdfstring{\ru{чайник} (cháynik)}{чайник (cháynik)}]{\ru{чайник} (cháynik) — a kettle}

\label{lesson:RU-C246-chaynik}



\begin{quote}
Before the new one: say the Russian for excellent, then the Russian for a microwave.
\end{quote}

\subsection*{You'll want to know: \ru{чайник}}
\textbf{\ru{чайник}} (\emph{cháynik}) — \textquotedblleft{}a kettle\textquotedblright{}.

\begin{grammarlens}[title={What you've built}]
One more word for this chapter.
\end{grammarlens}

\subsection*{Your turn}
\begin{itemize}
\raggedright
  \item \emph{Say it:} \emph{cháynik}
  \item \emph{Say it:} \emph{cháynik}, once more
  \item \emph{Say it:} say \emph{cháynik}
  \item \emph{Recall:} say the Russian for excellent, then the Russian for a microwave, then say \emph{cháynik} again
\end{itemize}

\begin{tcolorbox}[breakable,colback=teal!4,colframe=teal!35!black,title={Before you move on}]
What does \ru{чайник} mean? (\textquotedblleft{}A kettle\textquotedblright{}.) Say it once more.
\end{tcolorbox}
\section[\texorpdfstring{\ru{миска} (míska)}{миска (míska)}]{\ru{миска} (míska) — a bowl}

\label{lesson:RU-C246-miska}



\begin{quote}
Before the new one: say the Russian for a microwave, then the Russian for a kettle.
\end{quote}

\subsection*{You'll want to know: \ru{миска}}
\textbf{\ru{миска}} (\emph{míska}) — \textquotedblleft{}a bowl\textquotedblright{}.

\begin{grammarlens}[title={What you've built}]
One more word for this chapter.
\end{grammarlens}

\subsection*{Your turn}
\begin{itemize}
\raggedright
  \item \emph{Say it:} \emph{míska}
  \item \emph{Say it:} \emph{míska}, once more
  \item \emph{Say it:} say \emph{míska}
  \item \emph{Recall:} say the Russian for a microwave, then the Russian for a kettle, then say \emph{míska} again
\end{itemize}

\begin{tcolorbox}[breakable,colback=teal!4,colframe=teal!35!black,title={Before you move on}]
What does \ru{миска} mean? (\textquotedblleft{}A bowl\textquotedblright{}.) Say it once more.
\end{tcolorbox}
\section[\texorpdfstring{\ru{блюдце} (blyúdtse)}{блюдце (blyúdtse)}]{\ru{блюдце} (blyúdtse) — a saucer}

\label{lesson:RU-C246-blyudtse}



\begin{quote}
Before the new one: say the Russian for a kettle, then the Russian for a bowl.
\end{quote}

\subsection*{You'll want to know: \ru{блюдце}}
\textbf{\ru{блюдце}} (\emph{blyúdtse}) — \textquotedblleft{}a saucer\textquotedblright{}.

\begin{grammarlens}[title={What you've built}]
One more word for this chapter.
\end{grammarlens}

\subsection*{Your turn}
\begin{itemize}
\raggedright
  \item \emph{Say it:} \emph{blyúdtse}
  \item \emph{Say it:} \emph{blyúdtse}, once more
  \item \emph{Say it:} say \emph{blyúdtse}
  \item \emph{Recall:} say the Russian for a kettle, then the Russian for a bowl, then say \emph{blyúdtse} again
\end{itemize}

\begin{tcolorbox}[breakable,colback=teal!4,colframe=teal!35!black,title={Before you move on}]
What does \ru{блюдце} mean? (\textquotedblleft{}A saucer\textquotedblright{}.) Say it once more.
\end{tcolorbox}
\section[\texorpdfstring{\ru{кружка} (krúzhka)}{кружка (krúzhka)}]{\ru{кружка} (krúzhka) — a mug}

\label{lesson:RU-C246-kruzhka}



\begin{quote}
Before the new one: say the Russian for a bowl, then the Russian for a saucer.
\end{quote}

\subsection*{You'll want to know: \ru{кружка}}
\textbf{\ru{кружка}} (\emph{krúzhka}) — \textquotedblleft{}a mug\textquotedblright{}.

\begin{grammarlens}[title={What you've built}]
One more word for this chapter.
\end{grammarlens}

\subsection*{Your turn}
\begin{itemize}
\raggedright
  \item \emph{Say it:} \emph{krúzhka}
  \item \emph{Say it:} \emph{krúzhka}, once more
  \item \emph{Say it:} say \emph{krúzhka}
  \item \emph{Recall:} say the Russian for a bowl, then the Russian for a saucer, then say \emph{krúzhka} again
\end{itemize}

\begin{tcolorbox}[breakable,colback=teal!4,colframe=teal!35!black,title={Before you move on}]
What does \ru{кружка} mean? (\textquotedblleft{}A mug\textquotedblright{}.) Say it once more.
\end{tcolorbox}
